% PLACEHOLDER DATA in panel (b).
\begin{figure}[t]
\centering
\begin{minipage}[b]{0.43\linewidth}
\centering
\begin{tikzpicture}[scale=0.82, transform shape,
  font=\scriptsize,
  sv/.style={circle, draw=cCira, fill=cMem, inner sep=1.5pt, minimum size=5mm},
  lv/.style={rectangle, draw=cNoCue, fill=cNoCue!15, inner sep=1.5pt, minimum size=4.5mm},
  seen/.style={thick, cCira},
  held/.style={thick, dashed, cReset},
]
% connected design
\node[font=\scriptsize\bfseries] at (0.9,2.35) {connected};
\foreach \n/\y in {A/1.8,B/1.0,C/0.2} \node[sv] (s\n) at (0,\y) {\n};
\foreach \n/\y in {1/1.8,2/1.0,3/0.2} \node[lv] (l\n) at (1.8,\y) {\n};
\foreach \a/\b in {A/1,A/2,B/1,B/3,C/2,C/3} \draw[seen] (s\a) -- (l\b);
\foreach \a/\b in {A/3,B/2,C/1} \draw[held] (s\a) -- (l\b);
% disconnected control
\begin{scope}[xshift=2.75cm]
\node[font=\scriptsize\bfseries] at (0.9,2.35) {disconnected};
\foreach \n/\y in {A/1.8,B/1.0,C/0.2} \node[sv] (t\n) at (0,\y) {\n};
\foreach \n/\y in {1/1.8,2/1.0,3/0.2} \node[lv] (m\n) at (1.8,\y) {\n};
\foreach \a/\b in {A/1,A/2,B/1,B/2,C/3} \draw[seen] (t\a) -- (m\b);
\draw[held] (tA) -- (m3);
\end{scope}
\node[font=\tiny, align=left, anchor=north west] at (-0.3,-0.25) {\textcolor{cCira}{\rule[2pt]{8pt}{1pt}} seen combination \quad \textcolor{cReset}{\rule[2pt]{3pt}{1pt}\,\rule[2pt]{3pt}{1pt}} held out\\ circles: surfaces, squares: loads};
\end{tikzpicture}
\\[-2pt]{\scriptsize (a) observation graphs}
\end{minipage}\hfill
\begin{minipage}[b]{0.54\linewidth}
\centering
\begin{tikzpicture}
\begin{axis}[
  width=\linewidth, height=3.3cm,
  ybar, bar width=5.5pt, ymin=0, ymax=1.1,
  ylabel={held-out error / frozen}, label style={font=\scriptsize}, tick label style={font=\tiny},
  symbolic x coords={conn,disc}, xtick=data, xticklabels={connected, disconnected},
  enlarge x limits=0.45,
  legend style={font=\tiny, at={(0.5,1.02)}, anchor=south, legend columns=3, draw=none, fill=none},
]
\addplot[fill=cFrozen!50, draw=cFrozen] coordinates {(conn,0.92) (disc,0.93)}; \addlegendentry{independent}
\addplot[fill=cCont!50, draw=cCont] coordinates {(conn,0.66) (disc,0.74)}; \addlegendentry{continuous ctx.}
\addplot[fill=cNoCue!60, draw=cNoCue] coordinates {(conn,0.41) (disc,0.88)}; \addlegendentry{additive $\kappa_\times{=}0$}
\addplot[fill=cCira, draw=cCira] coordinates {(conn,0.33) (disc,0.84)}; \addlegendentry{\method (full kernel)}
\addplot[fill=cOracle!60, draw=cOracle] coordinates {(conn,0.24) (disc,0.26)}; \addlegendentry{oracle factors}
\node[anchor=north east, font=\tiny, text=red!70!black] at (rel axis cs:0.99,0.97) {placeholder};
\end{axis}
\end{tikzpicture}
\\[-2pt]{\scriptsize (b) first-ever-contact error on held-out cells}
\end{minipage}
\caption{\textbf{Composition before the first-ever contact} (placeholder data in b). (a)~Surfaces A--C and loads 1--3. In the connected design every held-out cell is linked to the seen cells by a path. In the disconnected control, the held-out cell A3 joins two components. (b)~Prediction error at the first contact of each held-out combination. The factor prior transfers only in the connected design, as \cref{prop:composition} predicts; the interaction term ($\kappa_\times>0$) helps when surface and load interact.}
\label{fig:composition}
\end{figure}
